% Level 3 workbook, Week 9 -- How Do Machines Learn?
\chapterweek{9}{How Do Machines Learn?}{AI finds patterns in examples -- and copies their mistakes.}

\section{What you will learn}
\begin{itemize}[itemsep=1pt]
  \item The difference between a program that follows \textbf{rules} and one that \textbf{learns from examples}.
  \item How a \textbf{nearest-neighbor classifier} sorts things, and how to measure its \textbf{accuracy}.
  \item Why unbalanced or wrong examples make an AI unfair (\textbf{bias}).
  \item How a tiny text generator writes sentences -- and why chatbots can be confidently wrong.
\end{itemize}

\section{Rules or examples?}
A program with \textbf{rules} does exactly what you typed: \py|if teeth > 5: print("grumpy")|. A \textbf{learning} program gets lots of \textbf{examples} with the right answers (\textbf{labels}) and finds the pattern itself. The things it measures are called \textbf{features}.

Our monsters have two features: \textbf{fluff} (0--10) and \textbf{teeth} (0--10). Put every example on a chart. To guess a \textbf{mystery} monster, find the closest example and copy its label. That is a \textbf{nearest-neighbor classifier}.

\begin{center}
\begin{tikzpicture}[x=0.5cm,y=0.5cm,>=Stealth]
  \draw[->] (0,0) -- (10.8,0) node[below left,font=\small] {fluff};
  \draw[->] (0,0) -- (0,10.8) node[above,font=\small] {teeth};
  \foreach \t in {0,5,10} {\node[below,font=\scriptsize] at (\t,0) {\t}; \node[left,font=\scriptsize] at (0,\t) {\t};}
  \draw[dashed,thick] (5,2) -- (6,1);
  \draw[dashed,thick] (3,6) -- (2,6);
  \foreach \x/\y in {9/1,8/2,7/3,6/1,8/4} {
    \draw[fill=kcorange!70] (\x,\y) circle (0.42); \node[font=\scriptsize\bfseries] at (\x,\y) {F};}
  \foreach \x/\y in {2/9,3/8,1/7,4/9,2/6} {
    \draw[fill=kcsky!70] (\x-0.4,\y-0.4) rectangle (\x+0.4,\y+0.4); \node[font=\scriptsize\bfseries] at (\x,\y) {G};}
  \foreach \x/\y in {5/2,3/6} {
    \draw[fill=white,thick] (\x,\y) circle (0.42); \node[font=\scriptsize\bfseries] at (\x,\y) {?};}
  \node[font=\small,align=left,anchor=west] at (11.5,8) {\tikz\draw[fill=kcorange!70] (0,0) circle (0.2); F = friendly (circle)\\[2pt]
    \tikz\draw[fill=kcsky!70] (0,0) rectangle (0.35,0.35); G = grumpy (square)\\[2pt]
    \tikz\draw[fill=white,thick] (0,0) circle (0.2); ? = mystery monster\\[2pt]
    dashed line = its nearest neighbor};
\end{tikzpicture}
\end{center}

The mystery at (5, 2) is closest to a friendly monster, so the guess is \emph{friendly}. Distance is measured with Pythagoras: \py|math.sqrt(dx * dx + dy * dy)|.

\section{Is it any good? Train, then test}
Split the examples into two piles. The program learns from the \textbf{training} pile. The \textbf{test} pile stays secret until the end; then we count how many test cards it labels correctly.
\[ \text{accuracy} = \frac{\text{correct guesses}}{\text{number of test cards}} \]
Testing on cards it learned from would be like seeing the exam answers first.

\begin{bughunt}[Bad examples, bad predictions]
If the training pile has 8 grumpy monsters and only 1 friendly one, the program calls most friendly monsters grumpy. The code is fine -- the \emph{data} is the bug. This is called \textbf{bias}. When real AI decides things about people (faces, voices, job applications), biased examples can treat whole groups unfairly. That is why people must check the data, not just the code.
\end{bughunt}

\section{A tiny text generator}
A \textbf{Markov chain} learns just one thing from a text: \emph{after each word, which words came next?}

\begin{center}
\begin{tikzpicture}[w/.style={draw,rounded corners=2pt,font=\small\ttfamily,fill=kcyellow!25,inner sep=3pt},>=Stealth]
  \node[w] (the) at (0,0) {the};
  \node[w] (dragon) at (3,1.2) {dragon};
  \node[w] (robot) at (3,0) {robot};
  \node[w] (cat) at (3,-1.2) {cat};
  \node[w] (sat) at (6,0) {sat};
  \node[w] (on) at (8.3,0) {on};
  \draw[->] (the) -- (dragon); \draw[->] (the) -- (robot); \draw[->] (the) -- (cat);
  \draw[->] (dragon) -- (sat); \draw[->] (robot) -- (sat); \draw[->] (cat) -- (sat);
  \draw[->] (sat) -- (on);
  \draw[->] (on) to[bend left=35] node[below,font=\scriptsize] {next word: the} (the);
\end{tikzpicture}
\end{center}
To write, it starts with a word and keeps picking a \emph{random} next word from the list. The sentences \emph{sound} like the training text, but the machine does not understand them at all. Today's chatbots are enormously bigger and cleverer, but they also predict likely next words -- which is why they can sound sure and still be wrong.

\section{The five big ideas of AI}
\begin{center}\small
\begin{tabular}{@{}p{3.4cm}p{6cm}p{4.4cm}@{}}
\toprule
\textbf{Big idea} & \textbf{What it means} & \textbf{Today}\\ \midrule
1 Perception & Computers sense the world with cameras, microphones, sensors. & a camera measuring fluff?\\
2 Representation \& reasoning & The world is stored as data the computer can reason about. & each monster = (fluff, teeth)\\
3 Learning & Computers find patterns in examples. & nearest-neighbor classifier\\
4 Natural interaction & Talking and writing like people is very hard. & Markov story machine\\
5 Societal impact & AI can help or harm people. & bias, privacy, deepfakes\\
\bottomrule
\end{tabular}
\end{center}

\begin{exercise}[Monster sorter]
Folder: \texttt{week09/ex-monster-sorter}. Finish \texttt{distance(a, b)} and \texttt{nearest(training, mystery)}. Run it and look at the chart. What does it guess for (5, 5)? Do you agree? \answerlines{1}
\end{exercise}

\begin{exercise}[Fair or biased?]
Folder: \texttt{week09/ex-fair-or-biased}. Finish \texttt{label\_counts()} and \texttt{accuracy()}.
Fair training accuracy: \rule{1.5cm}{0.4pt} \quad Biased training accuracy: \rule{1.5cm}{0.4pt}
Why is the biased one worse? \answerlines{1}
\end{exercise}

\begin{exercise}[Markov story machine]
Folder: \texttt{week09/ex-markov-story}. Finish \texttt{build\_chain(words)}: each word maps to the list of words that came right after it (keep repeats!). Copy your silliest generated sentence: \answerlines{1}
\end{exercise}

\begin{exercise}[My project pitch (Milestone 1)]
For the next three weeks you will \emph{Build Something Big}. Read the Project brief, then fill in your pitch. Finish it at home.
\begin{tabular}{@{}p{4.6cm}p{\dimexpr\linewidth-4.8cm-2\tabcolsep}@{}}
Project title: & \rule{\linewidth}{0.4pt}\\[6pt]
Option (circle one): & grid game \quad generative art app \quad data/AI app\\[6pt]
In one sentence, it is...: & \rule{\linewidth}{0.4pt}\\[6pt]
Two classes I will need: & \rule{\linewidth}{0.4pt}\\[6pt]
Inheritance or composition: & \rule{\linewidth}{0.4pt}\\[6pt]
Data structure (dict / stack / queue / grid): & \rule{\linewidth}{0.4pt}\\[6pt]
Named algorithm: & \rule{\linewidth}{0.4pt}\\[6pt]
File it will save: & \rule{\linewidth}{0.4pt}\\[6pt]
The smallest fun version: & \rule{\linewidth}{0.4pt}\\
\end{tabular}
\end{exercise}

\begin{safety}[Using AI wisely and safely]
\textbf{AI can be wrong}, even when it sounds sure -- check important facts with a trusted source or adult. \textbf{Never share personal information} with a chatbot: no real name, address, school, passwords, photos or secrets; what you type may be stored. \textbf{Deepfakes} are fake pictures, voices or videos made by AI: ask \enquote{Who made this, and why?} before you believe or share, and never make fakes of real people. \textbf{Bias}: AI is only as fair as its examples.
\end{safety}

\begin{vocab}
  \vocabword{feature}{Something we measure about an example, like fluff or teeth.}
  \vocabword{label}{The right answer for an example, like \enquote{friendly}.}
  \vocabword{training data}{The labeled examples a program learns from.}
  \vocabword{test data}{Examples kept secret during training, used to check the guesses.}
  \vocabword{accuracy}{Correct guesses divided by the number of test examples.}
  \vocabword{classifier}{A program that sorts things into groups (labels).}
  \vocabword{bias}{When unbalanced or unfair data makes a program treat some cases unfairly.}
  \vocabword{deepfake}{A fake photo, voice or video made with AI.}
\end{vocab}

\begin{reflection}
Where have you met AI this week (videos, games, phones, homework)? \answerlines{1}
One thing you will never type into a chatbot, and why: \answerlines{1}
\end{reflection}
